% ============================================================================
% preamble.tex -- all machinery, one edition (Polish explanations, German
% content).  The house colours are the estate's (research-documentation).
%
% CONTRACT: main.tex defines \bookstatus BEFORE it \inputs this file, because
% the running header reads it as it is read.
% ============================================================================
\usepackage[T1]{fontenc}
\usepackage[utf8]{inputenc}
\usepackage{lmodern}
\usepackage[ngerman,main=polish]{babel}
\usepackage[a4paper,left=2.4cm,right=2.4cm,top=2.6cm,bottom=2.6cm]{geometry}
\usepackage{microtype}
\usepackage{xcolor}
\usepackage{etoolbox}
\usepackage{enumitem}
\usepackage{booktabs}
\usepackage{longtable}
\usepackage{array}
\usepackage{colortbl}
\usepackage{multicol}
\usepackage{fancyhdr}
\usepackage{titlesec}
\usepackage[most]{tcolorbox}
\usepackage{amssymb}

% ---- house colours ---------------------------------------------------------
\definecolor{primaryblue}{RGB}{41,128,185}
\definecolor{darkblue}{RGB}{31,78,121}
\definecolor{lightgray}{RGB}{236,240,241}
\definecolor{darkgray}{RGB}{52,73,94}
\definecolor{accentgreen}{RGB}{39,174,96}
\definecolor{accentorange}{RGB}{230,126,34}
\definecolor{accentred}{RGB}{231,76,60}
% grammatical gender: the article is always printed, colour only reinforces it
\definecolor{colm}{RGB}{31,78,170}
\definecolor{colf}{RGB}{200,40,60}
\definecolor{coln}{RGB}{30,130,76}
\newcommand{\artm}[1]{\textcolor{colm}{\textbf{#1}}}
\newcommand{\artf}[1]{\textcolor{colf}{\textbf{#1}}}
\newcommand{\artn}[1]{\textcolor{coln}{\textbf{#1}}}

\setlength{\parindent}{0pt}
\setlength{\parskip}{0.55em plus 0.2em minus 0.1em}
\setlist{nosep,leftmargin=1.6em}
\renewcommand{\arraystretch}{1.12}

% ---- computed values: \val{key} reads generated/values.tex ----------------
% GENERATED by tools/gen.py -- do not edit.
\csdef{gval@candos.A1}{48}
\csdef{gval@candos.A2}{0}
\csdef{gval@candos.B1}{0}
\csdef{gval@candos.B2}{0}
\csdef{gval@candos.total}{48}
\csdef{gval@dialogs.A1}{13}
\csdef{gval@dialogs.A2}{0}
\csdef{gval@dialogs.B1}{0}
\csdef{gval@dialogs.B2}{0}
\csdef{gval@dialogs.total}{13}
\csdef{gval@exercises.A1}{154}
\csdef{gval@exercises.A2}{0}
\csdef{gval@exercises.B1}{0}
\csdef{gval@exercises.B2}{0}
\csdef{gval@exercises.total}{154}
\csdef{gval@nouns.das}{58}
\csdef{gval@nouns.der}{100}
\csdef{gval@nouns.die}{74}
\csdef{gval@swiss.A1}{11}
\csdef{gval@swiss.A2}{0}
\csdef{gval@swiss.B1}{0}
\csdef{gval@swiss.B2}{0}
\csdef{gval@swiss.total}{11}
\csdef{gval@traps.A1}{12}
\csdef{gval@traps.A2}{0}
\csdef{gval@traps.B1}{0}
\csdef{gval@traps.B2}{0}
\csdef{gval@traps.total}{12}
\csdef{gval@units.A1}{12}
\csdef{gval@units.A2}{0}
\csdef{gval@units.B1}{0}
\csdef{gval@units.B2}{0}
\csdef{gval@units.total}{12}
\csdef{gval@verify.A1}{3}
\csdef{gval@verify.A2}{0}
\csdef{gval@verify.B1}{0}
\csdef{gval@verify.B2}{0}
\csdef{gval@verify.total}{3}
\csdef{gval@vocab.A1}{519}
\csdef{gval@vocab.A2}{0}
\csdef{gval@vocab.B1}{0}
\csdef{gval@vocab.B2}{0}
\csdef{gval@vocab.total}{519}
\csdef{gval@vocab.verbs}{72}
\csdef{gval@weeks.total}{13}
\csdef{gval@work.A1}{12}
\csdef{gval@work.A2}{0}
\csdef{gval@work.B1}{0}
\csdef{gval@work.B2}{0}
\csdef{gval@work.total}{12}
\csdef{gval@written.A1}{12}
\csdef{gval@written.A2}{0}
\csdef{gval@written.B1}{0}
\csdef{gval@written.B2}{0}
\csdef{gval@written.total}{12}

% GENERATED by tools/gen.py -- do not edit.
\csdef{unittitle@01}{Dźwięki i powitania}
\csdef{unittitlede@01}{Laute und Begrüßungen}
\csdef{unitlevel@01}{A1}
\csdef{unitweek@01}{1}
\csdef{unitgrammar@01}{alfabet, wymowa, Sie/du}
\csdef{unitwords@01}{35}
\csdef{unittitle@02}{Przedstawianie się}
\csdef{unittitlede@02}{Sich vorstellen}
\csdef{unitlevel@02}{A1}
\csdef{unitweek@02}{2}
\csdef{unitgrammar@02}{sein, zaimki osobowe, pytania W-}
\csdef{unitwords@02}{51}
\csdef{unittitle@03}{Liczby i dane kontaktowe}
\csdef{unittitlede@03}{Zahlen und Kontaktdaten}
\csdef{unitlevel@03}{A1}
\csdef{unitweek@03}{3}
\csdef{unitgrammar@03}{czasowniki regularne, liczebniki 0-100}
\csdef{unitwords@03}{52}
\csdef{unittitle@04}{Rzeczowniki i rodzajniki}
\csdef{unittitlede@04}{Nomen und Artikel}
\csdef{unitlevel@04}{A1}
\csdef{unitweek@04}{4}
\csdef{unitgrammar@04}{rodzaj, liczba mnoga, ein/eine}
\csdef{unitwords@04}{37}
\csdef{unittitle@05}{Haben, biernik i przeczenie}
\csdef{unittitlede@05}{Haben, Akkusativ und Verneinung}
\csdef{unitlevel@05}{A1}
\csdef{unitweek@05}{5}
\csdef{unitgrammar@05}{haben, Akkusativ, kein/nicht}
\csdef{unitwords@05}{34}
\csdef{unittitle@06}{Rodzina i zawody}
\csdef{unittitlede@06}{Familie und Berufe}
\csdef{unitlevel@06}{A1}
\csdef{unitweek@06}{6}
\csdef{unitgrammar@06}{przymiotniki dzierżawcze}
\csdef{unitwords@06}{36}
\csdef{unittitle@07}{Czas i terminy}
\csdef{unittitlede@07}{Zeit und Termine}
\csdef{unitlevel@07}{A1}
\csdef{unitweek@07}{7}
\csdef{unitgrammar@07}{godziny, dni, miesiące, daty}
\csdef{unitwords@07}{54}
\csdef{unittitle@08}{Jedzenie i picie}
\csdef{unittitlede@08}{Essen und Trinken}
\csdef{unitlevel@08}{A1}
\csdef{unitweek@08}{8}
\csdef{unitgrammar@08}{möchte, restauracja, ceny}
\csdef{unitwords@08}{58}
\csdef{unittitle@09}{Zakupy i ilości}
\csdef{unittitlede@09}{Einkaufen und Mengen}
\csdef{unitlevel@09}{A1}
\csdef{unitweek@09}{9}
\csdef{unitgrammar@09}{czasowniki nieregularne (a→ä, e→i)}
\csdef{unitwords@09}{45}
\csdef{unittitle@10}{Mieszkanie i miejsca}
\csdef{unittitlede@10}{Wohnen und Orte}
\csdef{unitlevel@10}{A1}
\csdef{unitweek@10}{10}
\csdef{unitgrammar@10}{przyimki in/auf/an, Wo? Wohin?}
\csdef{unitwords@10}{55}
\csdef{unittitle@11}{Czasowniki modalne}
\csdef{unittitlede@11}{Modalverben}
\csdef{unitlevel@11}{A1}
\csdef{unitweek@11}{11}
\csdef{unitgrammar@11}{können, müssen, wollen, dürfen, Satzklammer}
\csdef{unitwords@11}{29}
\csdef{unittitle@12}{Plan dnia i tryb rozkazujący}
\csdef{unittitlede@12}{Tagesablauf und Imperativ}
\csdef{unitlevel@12}{A1}
\csdef{unitweek@12}{12}
\csdef{unitgrammar@12}{czasowniki rozdzielnie złożone, Imperativ}
\csdef{unitwords@12}{33}

\newcommand{\val}[1]{\ifcsname gval@#1\endcsname\csuse{gval@#1}\else\textbf{??}\fi}
\newcommand{\twodig}[1]{\ifnum#1<10 0\fi\number#1}

% ---- German and Polish text ------------------------------------------------
\newcommand{\de}[1]{\foreignlanguage{ngerman}{\textcolor{darkblue}{#1}}}
\newcommand{\deb}[1]{\foreignlanguage{ngerman}{\textcolor{darkblue}{\textbf{#1}}}}
\newcommand{\code}[1]{\texttt{#1}}
\newcommand{\ch}{\textit{CH}}
% a claim that needs a primary source or a native speaker before you rely on it
\newcommand{\verify}[1]{#1\textsuperscript{\textcolor{accentorange}{[?]}}}

% ---- language names inside babel ------------------------------------------
\addto\captionspolish{%
  \renewcommand{\chaptername}{Tydzień}%
  \renewcommand{\appendixname}{Dodatek}%
  \renewcommand{\contentsname}{Spis treści}}

% ---- headings ---------------------------------------------------------------
\titleformat{\part}[display]
  {\centering\color{primaryblue}\sffamily\bfseries\Huge}
  {}{0pt}{}
\titleformat{\chapter}[display]
  {\color{darkblue}\sffamily\bfseries}
  {\Large\color{primaryblue}\MakeUppercase{\chaptertitlename}\ \thechapter}{0.4em}{\Huge}[\vspace{0.3em}\titlerule]
\titlespacing*{\chapter}{0pt}{-10pt}{22pt}
\titleformat{\section}
  {\color{primaryblue}\Large\bfseries\sffamily}{\thesection}{0.6em}{}
\titleformat{\subsection}
  {\color{darkblue}\large\bfseries\sffamily}{}{0pt}{}
\titlespacing*{\subsection}{0pt}{1.1em}{0.3em}

% ---- running header: status rides on every page ----------------------------
\pagestyle{fancy}
\fancyhf{}
\renewcommand{\headrulewidth}{0.4pt}
\fancyhead[L]{\small\sffamily\color{darkgray}\nouppercase{\leftmark}}
\fancyhead[R]{\small\sffamily\color{accentred}\bookstatus}
\fancyfoot[C]{\small\thepage}
\fancypagestyle{plain}{\fancyhf{}\renewcommand{\headrulewidth}{0pt}\fancyfoot[C]{\small\thepage}}

% ---- boxes ------------------------------------------------------------------
\tcbset{mfa/.style={enhanced,breakable,boxrule=0pt,frame hidden,left=9pt,right=8pt,top=5pt,bottom=5pt,
  before skip=9pt,after skip=9pt,borderline west={3pt}{0pt}{#1},colback=#1!7}}
\newcommand{\boxlead}[2]{{\sffamily\bfseries\color{#1}#2.}\ }
\newtcolorbox{falle}[1][Pułapka dla Polaka]{mfa=accentred,before upper={\boxlead{accentred}{#1}}}
\newtcolorbox{szwajcaria}[1][Szwajcaria]{mfa=accentorange,before upper={\boxlead{accentorange}{#1}}}
\newtcolorbox{praca}[1][W pracy]{mfa=primaryblue,before upper={\boxlead{primaryblue}{#1}}}
\newtcolorbox{zasada}[1][Zasada]{mfa=accentgreen,before upper={\boxlead{accentgreen}{#1}}}
\newtcolorbox{uwaga}[1][Uwaga]{mfa=darkgray,before upper={\boxlead{darkgray}{#1}}}

% can-do goals: a checkbox list the learner ticks
\newenvironment{candos}
  {\begin{tcolorbox}[mfa=primaryblue]{\sffamily\bfseries\color{primaryblue}Po tym tygodniu potrafisz (\textit{Ich kann \ldots}):}\begin{itemize}[label=$\square$]}
  {\end{itemize}\end{tcolorbox}}

\newcommand{\weekend}[1]{%
  \begin{tcolorbox}[mfa=accentgreen]
  {\sffamily\bfseries\color{accentgreen}Zadanie weekendowe (produkcja, nie rozpoznawanie).}\ #1
  \end{tcolorbox}}

% ---- unit opener ------------------------------------------------------------
\newcommand{\unitopen}[1]{%
  \chapter[\csuse{unittitle@#1}]{\csuse{unittitle@#1}}%
  \begin{center}\small\sffamily\color{darkgray}
  \textbf{\csuse{unitlevel@#1}}\quad|\quad tydzień kalendarzowy \csuse{unitweek@#1}\quad|\quad
  \foreignlanguage{ngerman}{\csuse{unittitlede@#1}}\\
  gramatyka: \csuse{unitgrammar@#1}\quad|\quad nowe słowa: \csuse{unitwords@#1}
  \end{center}}

% ---- vocabulary: tables are GENERATED from data/vocab ------------------------
\newenvironment{vocabtab}
  {\begin{longtable}{@{}>{\raggedright\arraybackslash}p{0.38\linewidth}>{\raggedright\arraybackslash}p{0.30\linewidth}>{\raggedright\arraybackslash\color{darkgray}}p{0.27\linewidth}@{}}
   \toprule\textsf{\textbf{Deutsch}}&\textsf{\textbf{Polski}}&\textsf{\textbf{Uwagi}}\\\midrule\endhead
   \bottomrule\endlastfoot}
  {\end{longtable}}
\newcommand{\vrow}[3]{\de{#1} & #2 & #3\\}
\newcommand{\vocabtable}[1]{\input{generated/vocab-u#1}}
% the alphabetical index is paragraphs, not a table: longtable cannot live in multicols
\newenvironment{indextable}{\begin{multicols}{2}\small\raggedright}{\end{multicols}}
\newcommand{\irow}[3]{\par\noindent\hangindent=1em \de{#1} --- #2 \hfill{\color{darkgray}t.#3}\par}

% ---- calendar ----------------------------------------------------------------
\newenvironment{plantable}
  {\begin{longtable}{@{}rcrp{0.37\linewidth}p{0.32\linewidth}@{}}
   \toprule\textsf{\textbf{Tydz.}}&\textsf{\textbf{Poz.}}&\textsf{\textbf{Nr}}&\textsf{\textbf{Temat}}&\textsf{\textbf{Gramatyka}}\\\midrule\endhead
   \bottomrule\endlastfoot}
  {\end{longtable}}

% ---- dialogues: German left, Polish right, speaker label --------------------
\newenvironment{dialog}
  {\par\medskip\begin{longtable}{@{}>{\bfseries}p{1.5em}p{0.45\linewidth}>{\color{darkgray}\itshape}p{0.40\linewidth}@{}}}
  {\end{longtable}\par}
\newcommand{\turn}[3]{#1 & \de{#2} & #3\\}

% ---- exercises: answers are written where the exercise is, printed at the back
\newcommand{\answerstore}{}
\newcommand{\lastansunit}{0}
\newenvironment{cwiczenia}
  {\begin{enumerate}[label=\textbf{\arabic*.},leftmargin=2em,itemsep=0.35em]}
  {\end{enumerate}}
\newcommand{\exitem}[2]{%
  \item #1%
  \begingroup\edef\tmp{\endgroup\noexpand\gappto\noexpand\answerstore{%
    \noexpand\answeritem{\number\value{chapter}}{\number\value{enumi}}{\unexpanded{#2}}}}\tmp}
\newcommand{\answeritem}[3]{%
  \ifnum#1=\lastansunit\else
    \gdef\lastansunit{#1}\par\medskip\noindent{\sffamily\bfseries\color{primaryblue}Tydzień #1 --- \csname unittitle@\twodig{#1}\endcsname}\par\nobreak
  \fi
  \noindent\textbf{#1.#2}\ #3\par}

% ---- hyperref last -----------------------------------------------------------
\usepackage[colorlinks=true,linkcolor=primaryblue,urlcolor=primaryblue,citecolor=primaryblue,
  pdftitle={Deutsch von Null bis B2},pdfauthor={Konrad Cinkusz},pdflang={pl-PL},bookmarksnumbered=true]{hyperref}
